% =============================================================================
%  SECTION 3 — MATHEMATICAL FORMALISATION OF INTERNAL PHASE SYMMETRY
%  SECTION 4 — THE LAW OF AMPLITUDE AND STRUCTURAL REGIME ADJUSTMENTS
%  Integrate immediately after Section 2 in the main document.
% =============================================================================


% =============================================================================
\section{Mathematical Formalisation of Internal Phase Symmetry}
% =============================================================================

\subsection{Bitcoin Cycles as a Damped Harmonic System}

The empirical regularities documented in Section~2 motivate a formal analogical
framework borrowed from the theory of dynamic systems. We propose that the
observable structure of Bitcoin market cycles admits a tractable representation
as a \textit{discrete damped harmonic oscillator}: a system characterised by
near-constant periodicity, a stable internal phase relationship, and
monotonically decreasing amplitude across successive cycles.

In continuous-time mechanics, a damped harmonic oscillator obeys the
second-order linear differential equation:

\begin{equation}
  \ddot{x}(t) + 2\zeta\omega_0\,\dot{x}(t) + \omega_0^2\,x(t) = 0
  \label{eq:dho}
\end{equation}

where $x(t)$ denotes the displacement from equilibrium, $\omega_0$ the
natural angular frequency of the undamped system, and $\zeta \in (0,1)$ the
damping ratio. For $\zeta \in (0,1)$---the underdamped regime---the general
solution takes the form:

\begin{equation}
  x(t) = A_0\, e^{-\zeta\omega_0 t}
           \cos\!\bigl(\omega_d t + \varphi_0\bigr)
  \label{eq:dho_solution}
\end{equation}

where $\omega_d = \omega_0\sqrt{1 - \zeta^2}$ is the damped angular frequency,
$A_0$ is the initial amplitude, and $\varphi_0$ is the initial phase offset.
The solution describes oscillations of progressively decreasing amplitude
enveloped by the exponential decay factor $e^{-\zeta\omega_0 t}$.

We do not claim that equation~\eqref{eq:dho} governs Bitcoin's price path in
any literal mechanical sense. Markets are reflexive, non-linear systems, and
price dynamics are not solutions to differential equations with fixed
coefficients. The value of the analogy lies in its structural
correspondence: (i) a measurable and nearly constant period $T$; (ii) a
well-defined phase relationship pinned by an internal reference event; and
(iii) a monotonically decreasing amplitude envelope. All three properties are
empirically supported by the data in Section~2 for the temporal dimension and
are developed for the amplitude dimension in Section~4.

\subsection{Cycle Period and Angular Frequency}

The total cycle period $T_{\text{cycle}}$ is defined as the elapsed time
between two consecutive troughs:

\begin{equation}
  T_{\text{cycle}} = \Delta^{+}_i + \Delta^{-}_i
  \label{eq:period}
\end{equation}

Using the empirically derived averages $\bar{\Delta}^{+} \approx 1{,}060$ days
and $\bar{\Delta}^{-} \approx 370$ days:

\begin{equation}
  \bar{T}_{\text{cycle}} \approx 1{,}060 + 370 = 1{,}430 \text{ days}
  \label{eq:period_value}
\end{equation}

The corresponding mean angular frequency $\bar{\omega}$, expressed in
radians per day, is:

\begin{equation}
  \bar{\omega} = \frac{2\pi}{\bar{T}_{\text{cycle}}}
               = \frac{2\pi}{1{,}430}
               \approx 4.394 \times 10^{-3}\ \text{rad}\cdot\text{day}^{-1}
  \label{eq:omega}
\end{equation}

\subsection{Phase Definition and the Halving as a Temporal Node}

Within each cycle $\mathcal{C}_i = [T_i,\, P_i,\, T_{i+1}]$, we define the
instantaneous cycle phase $\phi(t)$ as the normalised position of any date $t$
within the expansion interval $[T_i, P_i]$:

\begin{equation}
  \phi(t) = \frac{t - T_i}{\Delta^{+}_i}, \quad \phi \in [0,\, 1]
  \label{eq:phase_norm}
\end{equation}

Under this normalisation, $\phi(T_i) = 0$ at the cycle trough and
$\phi(P_i) = 1$ at the cycle peak. The phase at the halving event $H_i$ is
therefore:

\begin{equation}
  \phi(H_i) = \frac{\alpha_i}{\Delta^{+}_i}
             = \frac{\alpha_i}{\alpha_i + \beta_i}
  \label{eq:phase_halving}
\end{equation}

Substituting the empirical averages $\bar{\alpha} = 523$ days and
$\bar{\beta} = 536$ days:

\begin{equation}
  \bar{\phi}(H) = \frac{523}{523 + 536}
                = \frac{523}{1{,}059}
                \approx 0.4939
  \label{eq:phase_halving_value}
\end{equation}

The halving consistently occurs at approximately $\bar{\phi}(H) \approx 0.494$
of the expansion phase---marginally before the mathematical midpoint
$\phi = 0.500$. We define the \textit{phase displacement} $\delta_\phi$ as the
deviation of the halving from perfect bisection:

\begin{equation}
  \delta_\phi = \bar{\phi}(H) - 0.500 = 0.494 - 0.500 = -0.006
  \label{eq:phase_displacement}
\end{equation}

This displacement of $-0.006$ in normalised phase corresponds to approximately
$-0.006 \times 1{,}059 \approx -6$ calendar days: the halving arrives,
on average, six days earlier than the geometric midpoint of the expansion. The
sign is negative because the system's peak lies slightly further in time from
the halving than the trough does---a structural asymmetry we interpret as the
lagged propagation of the supply shock through market sentiment and price
discovery mechanisms.

The observed inter-cycle values of $\phi(H_i)$ are:

\begin{align}
  \phi_1(H_1) &= \frac{542}{1068} \approx 0.508 \\[4pt]
  \phi_2(H_2) &= \frac{513}{1061} \approx 0.484 \\[4pt]
  \phi_3(H_3) &= \frac{515}{1050} \approx 0.490
\end{align}

The range of $\phi(H_i)$ is $[0.484,\; 0.508]$, a span of 0.024 in normalised
phase, or approximately 25 calendar days when projected onto the mean expansion
duration. This near-invariance of the halving's phase position across three
independent cycles---each occurring in markedly different macroeconomic and
structural environments---constitutes the strongest quantitative argument for
treating the halving as a \textit{temporal node}: a fixed point in the phase
space of the system, rather than an event whose timing relative to the cycle
is coincidental.

\begin{definition}[Temporal Node]
  The halving event $H_i$ is a \textit{temporal node} of cycle $\mathcal{C}_i$
  if and only if its normalised phase $\phi(H_i)$ is invariant across cycles
  up to a maximum deviation of $\varepsilon_\phi$, where $\varepsilon_\phi$ is
  small relative to unity. Empirically, $\varepsilon_\phi \approx 0.012$ (one
  half of the observed range), satisfying the condition $\varepsilon_\phi \ll 1$.
\end{definition}

\subsection{Implications for Forward Projection}

The phase-anchoring property of the halving has a direct and practically
consequential implication: knowing the date of the next halving $H_{i+1}$
and an estimate of the trough $T_{i+1}$, it is possible to construct a
probabilistic interval for the next peak $P_{i+1}$ without reference to any
price-level information. Specifically:

\begin{equation}
  \hat{P}_{i+1} = H_{i+1} + \bar{\beta} \pm \sigma_{\max}^{\beta}
                = H_{i+1} + 536 \pm 12 \ \text{ days}
  \label{eq:peak_projection}
\end{equation}

This formulation renders the temporal projection \textit{falsifiable} in a
strictly Popperian sense: it generates a dated prediction interval that can
be confirmed or refuted by subsequent market data, independently of any
assumptions about price magnitude. The projection for the current cycle is
developed in full in Section~6.


% =============================================================================
\section{The Law of Amplitude and Structural Regime Adjustments}
% =============================================================================

\subsection{Empirical Amplitude Data}

The second measurable degree of freedom in the harmonic model is the
\textit{amplitude} of each cycle: the magnitude of the price return from
trough to peak. We characterise this dimension through two complementary
empirical series. The first is the \textit{cycle multiplier} $\mathcal{M}_i$,
defined as the ratio of the peak price to the trough price within cycle $i$:

\begin{equation}
  \mathcal{M}_i = \frac{P^{\$}_i}{T^{\$}_i}
  \label{eq:multiplier}
\end{equation}

The second is the \textit{prior drawdown} $d_i$, defined as the proportional
decline from the preceding cycle peak $P^{\$}_{i-1}$ to the current trough
$T^{\$}_i$:

\begin{equation}
  d_i = \frac{P^{\$}_{i-1} - T^{\$}_i}{P^{\$}_{i-1}}
      = 1 - \frac{T^{\$}_i}{P^{\$}_{i-1}}, \quad d_i \in (0,\, 1)
  \label{eq:drawdown}
\end{equation}

Table~\ref{tab:amplitude} reports the empirical values of $\mathcal{M}_i$ and
$d_i$ for the three observed cycles.

\begin{table}[H]
  \centering
  \caption{Bitcoin Cycle Amplitude Parameters: Multipliers and Prior Drawdowns}
  \label{tab:amplitude}
  \renewcommand{\arraystretch}{1.35}
  \begin{tabularx}{\textwidth}{l S[table-format=2.2] S[table-format=1.2]
                                 S[table-format=2.1] c}
    \toprule
    \textbf{Cycle} &
    \multicolumn{1}{c}{$\mathcal{M}_i$ (Trough$\to$Peak)} &
    \multicolumn{1}{c}{$d_i$ (Prior drawdown)} &
    \multicolumn{1}{c}{Return (\%)} &
    \textbf{Regime} \\
    \midrule
    Cycle 1 (2015--2017) & 130.40 & 0.83 & {+12{,}940\%} & Retail \\
    Cycle 2 (2018--2021) &  20.44 & 0.59 & {+1{,}944\%}  & Retail \\
    Cycle 3 (2022--2025)\textsuperscript{\dag}
                          &   8.07 & 0.48 & {+707\%}      & Transitional \\
    \midrule
    \textit{Projected Cycle 4} & \textit{2.5--4.0} & \textit{0.30--0.40?}
      & \textit{+150--300\%} & \textit{Institutional} \\
    \bottomrule
  \end{tabularx}
  \medskip
  \begin{minipage}{\textwidth}
    \small
    \textsuperscript{\dag} Cycle 3 amplitude computed under the hypothesis
    $P_3 = 6$ October 2025. Cycle 4 projections are model-derived estimates
    subject to the regime-adjustment discussed in \S\ref{sec:etf_break}.
  \end{minipage}
\end{table}

Two secular trends are immediately apparent from Table~\ref{tab:amplitude}.
First, the cycle multiplier $\mathcal{M}_i$ declines sharply and
monotonically across successive cycles: from $130.4\times$ to $20.4\times$
to $8.1\times$. Second, the prior drawdown $d_i$ likewise declines: from $83\%$
to $59\%$ to $48\%$. The co-movement of these two series suggests that the
amplitude of each recovery is not independent of the severity of the preceding
contraction---a relationship that we formalise in the following subsection.

\subsection{The Power-Law of Amplitude: Empirical Derivation}

Inspection of the data in Table~\ref{tab:amplitude} suggests that the cycle
multiplier $\mathcal{M}_i$ is functionally related to the prior drawdown $d_i$
through a power-law specification of the form:

\begin{equation}
  \mathcal{M}(d) = \kappa \cdot d^{\,\gamma}
  \label{eq:amplitude_law}
\end{equation}

where $\kappa > 0$ is a scaling coefficient and $\gamma > 0$ is the elasticity
exponent. Fitting equation~\eqref{eq:amplitude_law} to the three observed
data points $\{(d_i, \mathcal{M}_i)\}_{i=1}^{3}$ via least-squares regression
in log-linear space---that is, regressing $\ln\mathcal{M}$ on $\ln d$---yields
the following empirical estimates:

\begin{equation}
  \boxed{\ \hat{\mathcal{M}}(d) \approx 328 \cdot d^{\,5.1}\ }
  \label{eq:amplitude_fit}
\end{equation}

with $\hat{\kappa} = 328$ and $\hat{\gamma} = 5.1$.

The high value of the exponent $\hat{\gamma} = 5.1$ reflects the extreme
non-linearity of the relationship: small reductions in drawdown depth produce
disproportionately large compressions in the subsequent multiplier. This is
the mathematical expression of the \textit{Law of Diminishing Returns} that
characterises maturing asset classes. Table~\ref{tab:amplitude_fit} reports
the in-sample fit of the model against observed values.

\begin{table}[H]
  \centering
  \caption{In-Sample Fit of the Power-Law Amplitude Model
           $\hat{\mathcal{M}}(d) = 328 \cdot d^{5.1}$}
  \label{tab:amplitude_fit}
  \renewcommand{\arraystretch}{1.35}
  \begin{tabularx}{\textwidth}{l
    S[table-format=1.2]
    S[table-format=3.1]
    S[table-format=3.1]
    S[table-format=2.1]}
    \toprule
    \textbf{Cycle} &
    \multicolumn{1}{c}{$d_i$} &
    \multicolumn{1}{c}{$\hat{\mathcal{M}}(d_i)$ (Model)} &
    \multicolumn{1}{c}{$\mathcal{M}_i$ (Observed)} &
    \multicolumn{1}{c}{Error (\%)} \\
    \midrule
    Cycle 1 & 0.83 & 126.0 & 130.4 & {-3.4\%} \\
    Cycle 2 & 0.59 &  22.0 &  20.4 & {+7.8\%} \\
    Cycle 3 & 0.48 &   7.7 &   8.1 & {-4.9\%} \\
    \bottomrule
  \end{tabularx}
\end{table}

The model reproduces the three observed multipliers with errors of less than
$8\%$, a precision that is considered high for a two-parameter empirical model
fitted to three data points across markets separated by a decade and orders of
magnitude in capitalisation. The direction of the errors is unsystematic
(negative, positive, negative), suggesting the absence of a structural bias in
the fit and consistent with random measurement or timing uncertainty rather
than model misspecification.

\begin{figure}[H]
    \centering
    \includegraphics[width=0.9\textwidth]{law_of_amplitude.png}
    \caption{\textbf{The Law of Amplitude and Regime Shift:} Empirical power-law trajectory mapped as a function of prior drawdown severity ($\hat{\mathcal{M}}(d) = 328 \cdot d^{5.1}$). The yellow boundary highlights the projected structural break and volatility dampening induced by the post-2024 spot institutional ETF regime change.}
    \label{fig:law_of_amplitude}
\end{figure}

\subsection{The Law of Diminishing Returns: Structural Interpretation}

The power-law relationship in equation~\eqref{eq:amplitude_fit} is not
merely a curve-fitting exercise. It admits a structural interpretation grounded
in the economics of asset maturation. We identify three reinforcing mechanisms
that produce the observed amplitude compression:

\textbf{Mechanism I: Market Capitalisation Scaling.} As Bitcoin's market
capitalisation increases by an order of magnitude with each cycle, the absolute
capital inflow required to achieve a given price multiplier grows
proportionally. A $100\times$ return from a \$1B market cap requires \$100B
in net new capital; the same return from a \$100B capitalisation requires
\$10T---a figure that exceeds the investable universe of most asset classes. The
\textit{reflexive compression} of multipliers as capitalisation grows is
therefore a near-mathematical inevitability rather than a contingent market
phenomenon.

\textbf{Mechanism II: Drawdown Compression as a Maturity Indicator.} The
monotonic decline in $d_i$ from $0.83$ to $0.59$ to $0.48$ is itself a
manifestation of structural maturation. As the asset develops deeper spot and
derivatives liquidity, a larger and more diversified investor base, and
greater macro-level awareness, the conditions for absolute capitulation---which
characterised the $-83\%$ drawdown of 2014--2015---become progressively less
tenable. Tighter drawdowns imply, via equation~\eqref{eq:amplitude_law} and
the large exponent $\hat{\gamma}=5.1$, dramatically lower subsequent
multipliers. In this sense, the drawdown depth $d_i$ functions as an implicit
\textit{maturity proxy} for the asset at the time of cycle trough.

\textbf{Mechanism III: Reflexive Expectation Dampening.} With each completed
cycle, the universe of potential Bitcoin investors who have not yet participated
shrinks. As the narrative of extraordinary returns becomes part of collective
financial memory, forward-looking participants price in a portion of the
expected cycle gain during the accumulation phase itself, reducing the
magnitude of the subsequent re-rating. This mechanism---analogous to the
well-documented reduction in the post-earnings announcement drift as analyst
coverage increases---produces a secular compression in realised multipliers
independently of capitalisation effects.

\subsection{The Structural Break: Spot ETF Approval as a Regime Shift}
\label{sec:etf_break}

The preceding mechanisms describe endogenous sources of amplitude compression
that would have been operative regardless of the regulatory and product
landscape. Beginning in January 2024, however, a qualitatively distinct and
exogenous force entered the system: the approval and large-scale adoption of
spot Bitcoin Exchange-Traded Funds in the United States and subsequently in
other major jurisdictions.

The ETF channel introduces capital flows with properties that are
fundamentally different from those of prior demand sources:

\begin{enumerate}
  \item \textbf{Passivity and mandate-driven allocation.} A substantial portion
  of ETF inflows originates from multi-asset portfolio rebalancing mandates
  rather than from discretionary conviction. This capital is, by construction,
  less sensitive to price levels and sentiment cycles than retail or
  speculative institutional capital.

  \item \textbf{Continuous and symmetric flow structure.} Unlike retail
  participants, who tend to concentrate purchases near cycle peaks and sales
  near troughs, mandate-driven ETF flows exhibit a more uniform distribution
  across time---dampening both the velocity of price appreciation during
  expansion phases and the velocity of decline during contraction phases.

  \item \textbf{Reduction of uninvested addressable market.} The ETF structure
  makes Bitcoin accessible to a class of capital---pension funds, endowments,
  wealth management platforms---that was previously excluded by custody,
  compliance, or fiduciary constraints. As this capital enters progressively
  during and after Cycle 3, the pool of ``yet-to-participate'' institutional
  demand diminishes, further constraining future multiplier potential.
\end{enumerate}

The net structural effect of ETF participation can be decomposed into two
directional impacts on the amplitude parameters:

\begin{equation}
  \Delta d_i^{\,\text{ETF}} < 0 \quad \Longrightarrow \quad
  \Delta \mathcal{M}_i^{\,\text{ETF}} < 0
  \label{eq:etf_effect}
\end{equation}

That is: ETF participation \textit{compresses the drawdown} $d_i$ below what
endogenous mechanisms would produce, which in turn---through the power-law
relationship of equation~\eqref{eq:amplitude_fit}---\textit{further compresses
the subsequent multiplier} $\mathcal{M}_i$. The two effects are not additive
but multiplicatively linked via the functional form of the amplitude law.

\subsection{Forward Amplitude Projection for Cycle 4}

Under the combined action of endogenous maturation and the ETF-driven
structural break, we project the prior drawdown for Cycle 4 in the range
$d_4 \in [0.30,\; 0.40]$. This range reflects:

\begin{itemize}
  \item The secular trend of declining drawdowns: $0.83 \to 0.59 \to 0.48$,
  suggesting a continuation toward the $[0.30, 0.40]$ interval.
  \item The additional stabilising force of sustained ETF demand during
  contraction phases, which introduces a structural price floor absent in
  prior cycles.
  \item Uncertainty about the magnitude of ETF-driven support, which precludes
  a point estimate and warrants the use of an interval.
\end{itemize}

Applying the power-law model of equation~\eqref{eq:amplitude_fit} to this
drawdown range, whilst acknowledging that the model was calibrated on
pre-ETF data and may therefore \textit{overestimate} the multiplier under the
new regime:

\begin{align}
  \hat{\mathcal{M}}(0.40) &= 328 \cdot (0.40)^{5.1}
    \approx 328 \cdot 0.00926 \approx 3.04 \label{eq:m_upper} \\[4pt]
  \hat{\mathcal{M}}(0.30) &= 328 \cdot (0.30)^{5.1}
    \approx 328 \cdot 0.00218 \approx 0.71 \label{eq:m_lower}
\end{align}

The lower bound produced by equation~\eqref{eq:m_lower} ($d_4 = 0.30$,
$\hat{\mathcal{M}} \approx 0.71$) implies a peak price below the trough price,
which is structurally inconsistent with a continuation of the harmonic cycle.
We interpret this as evidence that the purely mechanical extrapolation of the
power-law model breaks down at low drawdown levels---a regime where the
ETF-driven price floor prevents the kind of absolute capitulation that the
model was fitted to capture.

We therefore propose a \textit{regime-adjusted amplitude interval} for Cycle 4
that combines the mechanical model output for $d_4 \in [0.35, 0.40]$ with a
qualitative adjustment for ETF structural support:

\begin{equation}
  \mathcal{M}_4^{\text{adj}} \in [2.5,\; 4.0]
  \label{eq:m4_adjusted}
\end{equation}

This interval implies a percentage return from trough to peak of approximately
$+150\%$ to $+300\%$---a range that reflects the transition from a
hypergrowth asset to a maturing store-of-value instrument with deep
institutional participation. Critically, this projection is derived from the
structural and mechanical properties of the model rather than from price
extrapolation, and it should be treated as a probabilistic envelope rather
than a point forecast.

\begin{proposition}[Amplitude Convergence]
  Under continued asset maturation and sustained institutional ETF
  participation, the cycle multiplier $\mathcal{M}_i$ is expected to converge
  asymptotically toward a lower bound consistent with the return profile of
  other macro assets. The power-law relationship
  $\hat{\mathcal{M}}(d) = 328 \cdot d^{5.1}$ captures this convergence
  endogenously through the compression of prior drawdowns, but requires a
  regime adjustment for cycles in which ETF-driven demand materially
  alters the drawdown distribution relative to the pre-ETF calibration sample.
\end{proposition}

% =============================================================================
%  END OF SECTIONS 3 & 4
% =============================================================================
